\chapter{Architettura}

Il firmware segue un ciclo di esecuzione deterministico: inizializzazione in \texttt{setup()} e aggiornamento periodico dei moduli in \texttt{loop()}. Ogni modulo espone i metodi \texttt{begin()} e \texttt{update()} per separare configurazione e logica runtime, riducendo il rischio di stati inconsistenti.

\section{Schema logico}
\begin{itemize}
  \item \textbf{Turnstile}: legge i pulsanti IN/OUT, aggiorna il conteggio persone e muove il servo.
  \item \textbf{Stoplight}: visualizza su LED verde/rosso la disponibilita di posti.
  \item \textbf{Thermostat}: legge il sensore NTC e decide riscaldamento/raffreddamento.
  \item \textbf{LightingControl}: legge la fotoresistenza e regola il dimmer.
  \item \textbf{HumidifierControl}: legge il DHT22 e accende o spegne il controllo umidita.
  \item \textbf{DisplayPanel}: alterna pagine LCD con stato, data/ora e diagnostica.
\end{itemize}

\section{Flusso dati}
Il flusso segue tre passaggi costanti:
\begin{enumerate}
  \item \textbf{Acquisizione}: letture analogiche o digitali dai sensori.
  \item \textbf{Validazione}: controllo di errori e filtri sulle soglie.
  \item \textbf{Attuazione}: LED, PWM o servo in base alle decisioni.
\end{enumerate}
Il display aggrega gli stati dei moduli e presenta una vista sintetica per l'operatore.

\section{Cicli temporali}
Alcuni moduli usano tempi interni per evitare transizioni rapide o flicker. In particolare:
\begin{itemize}
  \item il termostato applica una pausa prima di cambiare modalita;
  \item il display limita il refresh a 1 secondo e cambia pagina ogni 5 secondi.
\end{itemize}
Queste scelte evitano oscillazioni e migliorano la stabilita percepita.

\section{Parametri principali}
I valori di target sono configurati nel file principale:\newline
\texttt{20.0 C} per temperatura, \texttt{65\%} per umidita, \texttt{200 lux} per illuminazione, \texttt{5 persone} come capienza massima.
